\clearpage\pagestyle{otMain}\fontsize{11}{15}\selectfont\corrigesChapitre{15}{Translations et vecteurs}
\begin{blocLivre}{cor-15-1}\begin{corrige}{1}{Symétrie et demi-tour}\demarche{Reporter les distances sur le quadrillage, point par point.}
Par la symétrie, chaque point est reporté perpendiculairement à (d), de l’autre côté et à la même distance : A et C sont à 6 carreaux de (d), B à 3 carreaux ; on obtient $\res{A'B'C'}$, tracé en vert. Par le demi-tour, O est le milieu de chaque segment point-image : on reproduit au-delà de O le déplacement de A vers O, soit 4 carreaux à droite et 3 en haut, pour placer A’’ ; de même pour $\res{B''\text{ et }C''}$. Les deux transformations $\res{\text{conservent les longueurs}}$.
\label{fig-cor-15-1}\figureGeo[corrige]{vect-01}
\end{corrige}\end{blocLivre}
\begin{blocLivre}{cor-15-2}\begin{corrige}{2}{Caractériser un vecteur}\demarche{Lire le déplacement horizontal et vertical, puis le reproduire.}
La direction est celle de la droite (AB), le sens va de A vers B. Le déplacement est de 4 carreaux vers la droite et 3 vers le haut : $AB=\sqrt{4^2+3^2}=\res{5\text{ cm}}$. On applique le même déplacement à partir de C : D est 4 carreaux à droite et 3 au-dessus de C.
\label{fig-cor-15-2}\figureGeo[corrige]{vect-02}
\end{corrige}\end{blocLivre}
\begin{blocLivre}{cor-15-3}\begin{corrige}{3}{Des vecteurs égaux}\demarche{Comparer direction, sens et longueur, ou le déplacement sur le quadrillage.}
$\vect{CD}$ et $\vect{IJ}$ correspondent au même déplacement que $\vect{AB}$ (2 à droite, 1 en haut) : $\res{\vect{CD}=\vect{IJ}=\vect{AB}}$. $\vect{EF}$ a même direction et même longueur, mais un sens contraire : $\res{\vect{EF}=-\vect{AB}=\vect{BA}}$. $\vect{GH}$ n’a pas la même direction.
\label{fig-cor-15-3}\figureGeo[corrige]{vect-03}
\end{corrige}\end{blocLivre}
\begin{blocLivre}{cor-15-4}\begin{corrige}{4}{Image par une translation}\demarche{Reproduire le déplacement, puis citer la définition.}
On applique à C le déplacement de A vers B : 3 carreaux à droite, 1 en haut. D’après la définition de la translation, $\res{ABDC\text{ est un parallélogramme}}$ et $\vect{CD}=\vect{AB}$. Les côtés [AB] et [CD] sont parallèles et de même longueur.
\label{fig-cor-15-4}\figureGeo[corrige]{vect-04}
\end{corrige}\end{blocLivre}
\begin{blocLivre}{cor-15-5}\begin{corrige}{5}{Vecteur nul et opposé}\demarche{Utiliser les définitions, puis la relation de Chasles.}
$\res{\vect{AA}=\vec0}$, le vecteur nul. L’opposé de $\vect{MN}$ est $\res{\vect{NM}}$ : même direction, même longueur, sens contraire. Par Chasles, $\vect{MN}+\vect{NM}=\vect{MM}=\res{\vec0}$.
\label{fig-cor-15-5}\figureGeo[corrige]{vect-05}
\end{corrige}\end{blocLivre}
\begin{blocLivre}{cor-15-6}\begin{corrige}{6}{Parallélogramme et égalités}\demarche{Suivre les côtés opposés dans le même sens.}
Dans ABCD, les côtés opposés donnent $\res{\vect{AB}=\vect{DC}}$ et $\res{\vect{AD}=\vect{BC}}$. En changeant le sens : $\res{\vect{BA}=\vect{CD}}$ et $\res{\vect{CB}=\vect{DA}}$. L’ordre des lettres suit le tour du parallélogramme.
\label{fig-cor-15-6}\figureGeo[corrige]{vect-06}
\end{corrige}\end{blocLivre}
\begin{blocLivre}{cor-15-7}\begin{corrige}{7}{Démontrer un parallélogramme}\demarche{Remplacer les lettres dans la propriété sans changer leur ordre.}
D’après la propriété du cours, $\vect{AB}=\vect{CD}$ équivaut à « ABDC est un parallélogramme ». Avec A = E, B = F, C = H et D = G, on obtient que $\res{EFGH\text{ est un parallélogramme}}$. Sur le quadrillage, les deux flèches correspondent au déplacement 4 à droite, 1 en haut.
\label{fig-cor-15-7}\figureGeo[corrige]{vect-07}
\end{corrige}\end{blocLivre}
\begin{blocLivre}{cor-15-8}\begin{corrige}{8}{La relation de Chasles}\demarche{Enchaîner les vecteurs bout à bout : l’extrémité de l’un est l’origine du suivant.}
$\res{\vect{AB}+\vect{BC}=\vect{AC}}$, comme sur la figure. $\vect{MN}+\vect{NP}+\vect{PQ}=\res{\vect{MQ}}$. En changeant l’ordre des termes, $\vect{CA}+\vect{AB}=\res{\vect{CB}}$. Enfin $\vect{AB}+\vect{BA}=\res{\vec0}$.
\label{fig-cor-15-8}\figureGeo[corrige]{vect-08}
\end{corrige}\end{blocLivre}
\begin{blocLivre}{cor-15-9}\begin{corrige}{9}{Construire une somme}\demarche{Mettre les représentants bout à bout.}
On trace $\vect{OP}=\vec u$ (3 à droite, 1 en haut), puis $\vect{PS}=\vec v$ (1 à droite, 2 en haut). Alors $\vec u+\vec v=\vect{OS}$, soit 4 à droite et 3 en haut : $OS=\sqrt{16+9}=\res{5\text{ cm}}$. Les longueurs de $\vec u$ et $\vec v$, $\sqrt{10}$ et $\sqrt5$, ont une somme différente.
\label{fig-cor-15-9}\figureGeo[corrige]{vect-09}
\end{corrige}\end{blocLivre}
\begin{blocLivre}{cor-15-10}\begin{corrige}{10}{Deux translations successives}\demarche{Remplacer chaque déplacement par un vecteur, puis additionner.}
$\vect{MM'}=\vect{AB}$ et $\vect{M'M''}=\vect{BC}$. Par Chasles, $\vect{MM''}=\vect{MM'}+\vect{M'M''}=\vect{AB}+\vect{BC}=\vect{AC}$ : c’est la $\res{\text{translation de vecteur }\vect{AC}}$. On le vérifie sur le quadrillage : 4 à droite, 3 en haut.
\label{fig-cor-15-10}\figureGeo[corrige]{vect-10}
\end{corrige}\end{blocLivre}
\begin{blocLivre}{cor-15-11}\begin{corrige}{11}{La règle du parallélogramme}\demarche{Remplacer un vecteur par un vecteur égal, puis appliquer Chasles.}
ABCD est un parallélogramme, donc $\vect{AD}=\vect{BC}$. Ainsi $\vect{AB}+\vect{AD}=\vect{AB}+\vect{BC}$, et la relation de Chasles donne $\res{\vect{AB}+\vect{AD}=\vect{AC}}$. La diagonale issue de A représente la somme.
\label{fig-cor-15-11}\figureGeo[corrige]{vect-11}
\end{corrige}\end{blocLivre}
\begin{blocLivre}{cor-15-12}\begin{corrige}{12}{Milieu et vecteurs}\demarche{Traduire « milieu » en direction, sens et longueur.}
I est sur [AB] à égale distance de A et B : $\vect{AI}$ et $\vect{IB}$ ont même direction, même sens et même longueur, donc $\res{\vect{AI}=\vect{IB}}$. Par Chasles, $\vect{AB}=\vect{AI}+\vect{IB}=\vect{AI}+\vect{AI}=\res{2\vect{AI}}$. Si $\vect{AJ}=\vect{JB}$, alors AJBJ est un parallélogramme aplati : ses diagonales [AB] et [JJ] ont même milieu, donc $\res{J\text{ est le milieu de }[AB]}$.
\label{fig-cor-15-12}\figureGeo[corrige]{vect-12}
\end{corrige}\end{blocLivre}
\begin{blocLivre}{cor-15-13}\begin{corrige}{13}{Un drone}\demarche{Additionner les déplacements, puis calculer la longueur avec Pythagore.}
Le premier déplacement est $\vect{DE}$, le second $\vect{EF}$. Par Chasles, leur somme est $\res{\vect{DF}}$. Le triangle DEF est rectangle en E : $DF=\sqrt{60^2+80^2}=\res{100\text{ m}}$, alors que le drone a parcouru 140 m.
\label{fig-cor-15-13}\figureGeo[corrige]{vect-13}
\end{corrige}\end{blocLivre}
\begin{blocLivre}{cor-15-14}\begin{corrige}{14}{Une frise}\demarche{Compter les intervalles entre les motifs.}
Du motif 1 au motif 6, il y a cinq déplacements : $\vec u+\vec u+\vec u+\vec u+\vec u$, noté $5\vec u$, de longueur $\res{20\text{ cm}}$. Du motif 6 au motif 2, on revient de quatre motifs : $\res{-4\vec u}$, même direction, sens contraire, longueur $\res{16\text{ cm}}$.
\label{fig-cor-15-14}\figureGeo[corrige]{vect-frise}
\end{corrige}\end{blocLivre}
\begin{blocLivre}{cor-15-15}\begin{corrige}{15}{Démontrer avec des vecteurs}\demarche{Chaîner deux égalités de vecteurs, puis citer la propriété.}
ABCD est un parallélogramme, donc $\vect{DC}=\vect{AB}$. Par hypothèse $\vect{BE}=\vect{AB}$. Donc $\vect{DC}=\vect{BE}$, ce qui signifie que $\res{DCEB\text{ est un parallélogramme}}$, d’après la propriété du cours appliquée aux points D, C, B, E.
\label{fig-cor-15-15}\figureGeo[corrige]{vect-15}
\end{corrige}\end{blocLivre}
